\documentclass[12pt,a4paper]{article}

% Packages
\usepackage[utf8]{inputenc}
\usepackage[bahasa]{babel}  % Indonesian language support
\usepackage{amsmath}
\usepackage{amsfonts}
\usepackage{amssymb}
\usepackage{graphicx}
\usepackage{hyperref}
\usepackage{geometry}
\geometry{margin=1in}

% Title setup
\title{Mencapai Optimisasi Portfolio Berkualitas Tinggi dengan VQE\\
\large Fase 0: Penurunan Teoritis}
\author{Tim DreamyTA}
\date{\today}

\begin{document}

\maketitle

\tableofcontents
\newpage

\section{Pemetaan QUBO ke Hamiltonian Ising}

\subsection{Formulasi Fungsi Objektif}

Masalah optimisasi portfolio dirumuskan sebagai masalah QUBO:
\begin{equation}
    C(x) = -\lambda \sum_{i=1}^{N} \mu_i x_i + (1-\lambda) \sum_{i=1}^{N} \sum_{j=i+1}^{N} \sigma_{ij} x_i x_j + p \left( \sum_{i=1}^{N} x_i - \frac{N}{2} \right)^2
\end{equation}

Di mana:
\begin{itemize}
    \item $x_i \in \{0, 1\}$: Variabel keputusan biner.
    \item $N$: Jumlah total aset.
    \item $\mu_i$: Pengembalian yang diharapkan dari aset $i$.
    \item $\sigma_{ij}$: Kovarian antara aset $i$ dan $j$.
    \item $\lambda$: Parameter penghindaran risiko ($0 \le \lambda \le 1$).
    \item $p$: Koefisien penalti untuk batasan kardinalitas ($\sum x_i = N/2$).
\end{itemize}

\subsection{Ekspansi Batasan}

Mengekspansi suku penalti kuadrat:
\begin{align*}
    \left( \sum_{i=1}^{N} x_i - \frac{N}{2} \right)^2 &= \left( \sum_{i=1}^{N} x_i \right)^2 - N \sum_{i=1}^{N} x_i + \frac{N^2}{4} \\
    &= \sum_{i=1}^{N} x_i + 2 \sum_{i<j} x_i x_j - N \sum_{i=1}^{N} x_i + \text{konst} \\
    &= \sum_{i=1}^{N} (1 - N) x_i + 2 \sum_{i<j} x_i x_j + \text{konst}
\end{align*}

Mensubstitusi kembali ke dalam $C(x)$ menghasilkan koefisien yang dikelompokkan.

\subsection{Pengelompokan Koefisien}

Menulis ulang $C(x)$ dengan koefisien linier ($Q_i$) dan kuadratik ($Q_{ij}$):
\begin{equation}
    C(x) = \sum_{i=1}^{N} Q_i x_i + \sum_{i=1}^{N} \sum_{j=i+1}^{N} Q_{ij} x_i x_j + \text{konst}
\end{equation}

\textbf{Koefisien:}
\begin{itemize}
    \item \textbf{Linier ($Q_i$):} $Q_i = -\lambda \mu_i + p(1 - N)$
    \item \textbf{Kuadratik ($Q_{ij}$):} $Q_{ij} = (1-\lambda) \sigma_{ij} + 2p$
\end{itemize}

\subsection{Transformasi ke Hamiltonian Ising}

Pemetaan ke model Ising menggunakan $x_i = \frac{1 - Z_i}{2}$, dengan $Z_i \in \{1, -1\}$ (nilai eigen Pauli-Z).

\textbf{1. Suku Linier:}
$$ \sum_{i} Q_i x_i = \sum_{i} \frac{Q_i}{2} - \sum_{i} \frac{Q_i}{2} Z_i $$

\textbf{2. Suku Kuadratik:}
$$ \sum_{i<j} Q_{ij} x_i x_j = \sum_{i<j} \frac{Q_{ij}}{4} (1 - Z_i - Z_j + Z_i Z_j) $$

\subsection{Ekspresi Hamiltonian Akhir}

Hamiltonian diberikan oleh $H = \sum_{i} h_i Z_i + \sum_{i<j} J_{ij} Z_i Z_j$, dengan koefisien:

\begin{equation}
    h_i = \frac{\lambda \mu_i}{2} - \frac{1-\lambda}{4} \sum_{j \neq i}^N \sigma_{ij}
\end{equation}

\begin{equation}
    J_{ij} = \frac{(1-\lambda) \sigma_{ij} + 2p}{4}
\end{equation}

\section{VQE dan CVaR}

\subsection{Latar Belakang VQE dan CVaR}

\textbf{Tujuan VQE:} Meminimalkan $E(\boldsymbol{\theta}) = \langle \psi(\boldsymbol{\theta}) | \hat{H} | \psi(\boldsymbol{\theta}) \rangle$.

\textbf{CVaR Standar (Conditional Value-at-Risk):} Alih-alih rata-rata, gunakan rata-rata dari energi ekor. Untuk energi yang diurutkan $E_{(1)} \le \dots \le E_{(K)}$:

\begin{equation}
    CVaR_\alpha(E) = \frac{1}{\lceil \alpha K \rceil} \sum_{k=1}^{\lceil \alpha K \rceil} E_{(k)}
\end{equation}

Ketika $\alpha = 1$, ini menghasilkan energi rata-rata standar.

\section{Weighted CVaR (WCVaR)}

\subsection{Definisi}

\begin{equation}
    WCVaR_\alpha(E) = \sum_{k=1}^{\lceil \alpha K \rceil} w_k E_{(k)}
\end{equation}
dengan batasan $\sum w_k = 1$.

\textbf{Perbedaan dari CVaR Standar:}
\begin{itemize}
    \item \textbf{CVaR Standar:} Bobot seragam $w_k = \frac{1}{\lceil \alpha K \rceil}$.
    \item \textbf{WCVaR:} Bobot tidak seragam, biasanya meluruh secara eksponensial.
\end{itemize}

\subsection{Manfaat Teoritis WCVaR}

\begin{itemize}
    \item \textbf{Lanskap yang Lebih Halus:} Tidak seperti optimisasi energi minimum ($\alpha \to 0$), WCVaR memberikan lanskap yang lebih halus dan dapat diturunkan.
    \item \textbf{Menghindari Dataran Tandus:} Tidak seperti ekspektasi penuh ($\alpha=1$), ini berfokus pada keadaan energi rendah, mencegah derau energi tinggi menghapus gradien.
    \item \textbf{Panduan:} Bobot eksponensial mendorong pengoptimal secara agresif menuju keadaan dasar sejati ($E_{(1)}$).
\end{itemize}

\subsection{Skema Pembobotan}

Beberapa skema pembobotan yang dapat digunakan (sebagaimana dijelaskan dalam Lampiran A):

\begin{enumerate}
    \item \textbf{Berbasis Energi:} $w_k = \exp [ -\beta(E_{(k)} - E_0) ]$ (Performa buruk).
    \item \textbf{Berbasis Peringkat:} $w_k = \exp(-\beta k)$ (Sederhana dan efektif).
    \item \textbf{Eksponensial Bertahap (Dipilih):}
    $$ w_k = \begin{cases} 
        \exp(-\beta_1 k) & k < N_1 \\
        \exp(-\beta_2 (k - N_1)) \cdot w_{N_1 - 1} & N_1 \le k < N_2 \\
        \dots & \dots
    \end{cases} $$
\end{enumerate}

Metode bertahap memungkinkan kontrol yang lebih rinci dan digunakan untuk hasil utama makalah.

\section{Optimisasi CMA-ES}

\subsection{Gambaran Umum}

\textbf{Covariance Matrix Adaptation Evolution Strategy (CMA-ES)} adalah algoritma optimisasi numerik stokastik tanpa turunan yang kuat dengan karakteristik berikut:

\begin{itemize}
    \item Algoritma optimisasi numerik stokastik tanpa turunan yang kuat.
    \item \textbf{Tipe:} Algoritma Evolusioner (domain kontinu).
    \item \textbf{Kasus Penggunaan:} Fungsi objektif non-konveks, tidak terkondisi dengan baik, multimodal, atau berisik.
    \item \textbf{Relevansi dengan VQE:} Digunakan untuk mengoptimalkan parameter $\boldsymbol{\theta}$ dari sirkuit variasional ketika gradien sulit dihitung atau berisik (misalnya, pada perangkat keras NISQ).
\end{itemize}

\subsection{Cara Kerja CMA-ES}

Algoritma mempertahankan distribusi Gaussian $\mathcal{N}(m, \sigma^2 C)$ dan memperbaruinya secara iteratif:

\begin{enumerate}
    \item \textbf{Pengambilan Sampel:} Menghasilkan populasi solusi kandidat baru (keturunan) dari distribusi saat ini.
    \item \textbf{Seleksi:} Mengevaluasi kebugaran (energi) dan memilih kandidat terbaik.
    \item \textbf{Adaptasi:} Memperbarui parameter distribusi:
        \begin{itemize}
            \item \textbf{Mean ($m$):} Bergerak menuju solusi yang lebih baik.
            \item \textbf{Ukuran Langkah ($\sigma$):} Mengontrol eksplorasi vs. eksploitasi.
            \item \textbf{Matriks Kovarian ($C$):} Mempelajari bentuk lanskap (korelasi antar variabel) untuk memandu pencarian.
        \end{itemize}
\end{enumerate}

\subsection{Ilustrasi Proses CMA-ES}

Proses optimisasi CMA-ES dapat divisualisasikan melalui beberapa langkah:

\begin{enumerate}
    \item \textbf{Inisialisasi:} Pencarian dimulai dengan mean awal dan distribusi yang luas.
    \item \textbf{Pengambilan Sampel \& Seleksi:} Kandidat diambil sampelnya; yang terbaik (biru) dipilih untuk memandu pembaruan.
    \item \textbf{Adaptasi:} Distribusi bergeser dan membentuk ulang (pembaruan kovarian) menuju lembah.
    \item \textbf{Eksplorasi:} Algoritma mengeksplorasi lanskap, memanjang sepanjang arah optimal.
    \item \textbf{Konvergensi Dimulai:} Ukuran langkah menurun saat populasi berpusat pada minimum.
    \item \textbf{Penyetelan Halus:} Lokalisasi presisi dari optimum.
    \item \textbf{Optimisasi Selesai:} Distribusi konvergen ke minimum global.
\end{enumerate}

\section{Kesimpulan}

Dokumen ini menyajikan penurunan teoritis untuk optimisasi portfolio menggunakan VQE dengan pendekatan WCVaR dan CMA-ES. Kombinasi metode ini memberikan kerangka kerja yang kuat untuk menyelesaikan masalah optimisasi portfolio yang kompleks pada perangkat keras kuantum NISQ.

\end{document}
